\subsection{RQ2: geometry, mapping and useful work}
\label{sec:rq2}

Fig.~\ref{fig:tradeoff} separates three quantities for the five 16-PE
shapes, using the same problem sizes as Table~\ref{tab:measured} and counted
array steps. The useful fraction $\eta$ is $1/16$ for vector addition, dot
product and scan on every shape: the diagonal produces $L$ results in $L$
steps on $RC$ PEs, the dot product uses one PE per step, and a scan block of
$C$ elements occupies the $C$ PEs of row~0 for $C$ steps. Matrix--vector
reaches $\eta=1/R$ (all sixteen PEs multiply on $1\times16$, one on
$16\times1$), and convolution less because band tiles contain structural
zeros. Static occupancy (panel~b) is up to sixteen times larger than $\eta$
for scan on $1\times16$, where all sixteen PEs are configured but only one
result per step is new, and should not be used as a throughput proxy.

For vector, dot and scan the shape changes only communication. The same 256-element vector
addition needs 2\,115 bytes on $4\times4$ but 11\,331 bytes on the strips,
because every transfer carries $R+C$ edge words for at most $\min(R,C)$
results; scan and matrix--vector are cheapest on $1\times16$ (1\,251 and
4\,870 bytes), and matrix--vector and convolution are most expensive on
$16\times1$ (47\,200 and 60\,608 bytes), where they also have the lowest
$\eta$. The number of steps of
matrix--vector ranges from 64 to 1\,024 across shapes, but, as shown for
RQ1, steps hardly matter on this link. No shape minimises traffic for every
kernel; for a known dominant kernel the shape should follow its edge
utilisation.

\begin{figure*}[t]
  \centering
  \resizebox{\textwidth}{!}{% Five 16-PE shapes. (a) useful fraction eta (Eq. eta) from counted steps;
% (b) static occupancy of arithmetic-configured PEs (model); (c) counted v3
% wire bytes per call. Data: build/ana-geometry.csv (tools/analysis_data.py).
\begin{tikzpicture}
  \tikzset{every axis/.append style={ybar, width=0.37\textwidth, height=44mm,
    symbolic x coords={1x16,2x8,4x4,8x2,16x1}, xtick=data,
    xticklabels={$1{\times}16$,$2{\times}8$,$4{\times}4$,$8{\times}2$,$16{\times}1$},
    tick label style={font=\scriptsize}, label style={font=\scriptsize},
    title style={font=\scriptsize, yshift=-1ex},
    axis line style={cWire!45}, tick style={cWire!45},
    ymajorgrids, grid style={cWire!15}, enlarge x limits=0.13, bar width=1.35mm}}
  \begin{axis}[name=a, title={(a) useful fraction $\eta$ (counted steps)},
      ymin=0, ymax=1.05, ylabel={$\eta$},
      legend style={font=\tiny, draw=cWire!35, fill=white, legend columns=-1,
                    at={(1.6,1.17)}, anchor=south}]
    \addplot[draw=cBlockE, fill=cBlock] table[x=shape, y=vector_useful, col sep=comma] {build/ana-geometry.csv};
    \addlegendentry{vector}
    \addplot[draw=cStateE, fill=cState] table[x=shape, y=dot_useful, col sep=comma] {build/ana-geometry.csv};
    \addlegendentry{dot}
    \addplot[draw=cRegE, fill=cReg] table[x=shape, y=scan_useful, col sep=comma] {build/ana-geometry.csv};
    \addlegendentry{scan}
    \addplot[draw=cAluE, fill=cAlu] table[x=shape, y=matvec_useful, col sep=comma] {build/ana-geometry.csv};
    \addlegendentry{matvec}
    \addplot[draw=cWire, fill=cWire!25] table[x=shape, y=conv_useful, col sep=comma] {build/ana-geometry.csv};
    \addlegendentry{conv}
  \end{axis}
  \begin{axis}[name=b, at={(a.east)}, anchor=west, xshift=11mm,
      title={(b) static arithmetic occupancy}, ymin=0, ymax=1.05, ylabel={fraction of PEs}]
    \addplot[draw=cBlockE, fill=cBlock] table[x=shape, y=vector_static, col sep=comma] {build/ana-geometry.csv};
    \addplot[draw=cStateE, fill=cState] table[x=shape, y=dot_static, col sep=comma] {build/ana-geometry.csv};
    \addplot[draw=cRegE, fill=cReg] table[x=shape, y=scan_static, col sep=comma] {build/ana-geometry.csv};
    \addplot[draw=cAluE, fill=cAlu] table[x=shape, y=matvec_static, col sep=comma] {build/ana-geometry.csv};
    \addplot[draw=cWire, fill=cWire!25] table[x=shape, y=conv_static, col sep=comma] {build/ana-geometry.csv};
  \end{axis}
  \begin{axis}[at={(b.east)}, anchor=west, xshift=11mm,
      title={(c) wire bytes per call, v3 (counted)}, ymode=log, log origin=infty,
      ymin=500, ymax=1e5, ylabel={bytes}]
    \addplot[draw=cBlockE, fill=cBlock] table[x=shape, y=vector_bytes, col sep=comma] {build/ana-geometry.csv};
    \addplot[draw=cStateE, fill=cState] table[x=shape, y=dot_bytes, col sep=comma] {build/ana-geometry.csv};
    \addplot[draw=cRegE, fill=cReg] table[x=shape, y=scan_bytes, col sep=comma] {build/ana-geometry.csv};
    \addplot[draw=cAluE, fill=cAlu] table[x=shape, y=matvec_bytes, col sep=comma] {build/ana-geometry.csv};
    \addplot[draw=cWire, fill=cWire!25] table[x=shape, y=conv_bytes, col sep=comma] {build/ana-geometry.csv};
  \end{axis}
\end{tikzpicture}
}
  \caption{Geometry and mapping for the five 16-PE shapes (workloads of
  Table~\ref{tab:measured}). (a)~Useful fraction $\eta$ of PE slots,
  Eq.~\eqref{eq:eta}, from counted steps: vector, dot and scan stay at
  $1/16$ on every shape; matrix--vector reaches $1/R$. (b)~Static fraction of
  PEs configured with an arithmetic opcode, which overstates efficiency
  wherever PEs recompute settled values or forward. (c)~Counted wire bytes
  of one v3 call. For vector, dot and scan the shape changes only
  communication; for the matrix kernels it changes both.}
  \label{fig:tradeoff}
\end{figure*}

